% =====================================================================
%  CARNET D'ENTRAINEMENT — FICHE 6 — THÈME 1
%  Niveau DIFFICILE — Corrigé
% =====================================================================
\documentclass[11pt,a4paper]{article}
\usepackage[utf8]{inputenc}
\usepackage[T1]{fontenc}
\usepackage{lmodern}
\usepackage[french]{babel}
\usepackage{amsmath,amssymb,mathtools}
\usepackage{icomma}
\usepackage[version=4]{mhchem}
\usepackage{siunitx}
\sisetup{output-decimal-marker={,}}
\usepackage{xcolor}
\usepackage{colortbl}
\usepackage{tikz}
\usepackage[most]{tcolorbox}
\usepackage{enumitem}
\usepackage{array}
\usepackage{booktabs}
\usepackage{fancyhdr}
\usepackage[hidelinks]{hyperref}
\usetikzlibrary{arrows.meta,positioning,calc,shapes.geometric}
\usepackage[a4paper,margin=2.1cm,headheight=26pt,headsep=8pt]{geometry}

\definecolor{primary}{RGB}{146,95,160}
\definecolor{primarydark}{RGB}{92,52,104}
\definecolor{corcol}{RGB}{0,114,178}
\definecolor{astucecol}{RGB}{198,93,9}

\tcbset{boxbase/.style={enhanced,breakable,boxrule=0.9pt,arc=2.5pt,
   left=3mm,right=3mm,top=2.2mm,bottom=2.2mm,
   fonttitle=\bfseries,coltitle=white,toptitle=0.6mm,bottomtitle=0.6mm}}
\newtcolorbox[auto counter]{cor}[1][]{boxbase,colback=corcol!5,colframe=corcol,
   title={\textsc{Corrigé}~\thetcbcounter\ifx\relax#1\relax\relax\else\textmd{ —\itshape\ #1}\fi}}
\newtcolorbox{astucebox}[1][]{boxbase,colback=astucecol!8,colframe=astucecol,
   title={\textsc{À retenir}\ifx\relax#1\relax\relax\else\textmd{ : \itshape #1}\fi}}

\pagestyle{fancy}\fancyhf{}
\fancyhead[L]{\small\textcolor{primarydark}{\textbf{Fiche 6 · Thème 1} — Corrigé}}
\fancyhead[R]{\small\textcolor{primarydark}{\textsc{Difficile}}}
\fancyfoot[C]{\small\thepage}
\renewcommand{\headrulewidth}{0.8pt}
\renewcommand{\headrule}{\hbox to\headwidth{\color{primary}\leaders\hrule height \headrulewidth\hfill}}
\setlength{\parindent}{0pt}\setlength{\parskip}{3pt}
\newcommand{\rep}{\textbf{\textcolor{corcol}{Réponse : }}}

\begin{document}
\begin{center}
{\Large\bfseries\color{primarydark}Thème 1 — Corrigé du niveau Difficile}
\end{center}
\medskip

\begin{cor}[la formation du sel de table, pas à pas]
\begin{enumerate}[label=\alph*),leftmargin=2em,itemsep=2pt]
  \item \ce{Na} : 1 électron de valence (colonne 1). \ce{Cl} : 7 électrons de valence (colonne 17).
  \item \ce{Na} n'a qu'\emph{un} électron externe : il a intérêt à le \textbf{céder}. \ce{Cl}, à
        qui il manque \emph{un} électron pour l'octet, a intérêt à le \textbf{capter}. On obtient
        \[ \ce{Na -> Na+ + e-}\qquad\text{et}\qquad \ce{Cl + e- -> Cl-}. \]
  \item \ce{Na+} a perdu son unique électron externe : sa nouvelle couche externe est complète,
        c'est l'\textbf{octet} du néon \ce{Ne}. \ce{Cl-} a gagné un électron : il atteint
        l'\textbf{octet} de l'argon \ce{Ar}. \emph{(Aucun des deux n'est un duet : le duet ne
        concerne que \ce{H} et \ce{He}.)}
  \item C'est une \textbf{liaison ionique}, assurée par l'\textbf{attraction électrostatique}
        (force de Coulomb) entre le cation \ce{Na+} et l'anion \ce{Cl-}, de charges opposées.
\end{enumerate}
\end{cor}

\begin{cor}[l'ion ammonium et la liaison dative]
\begin{enumerate}[label=\alph*),leftmargin=2em,itemsep=2pt]
  \item L'azote a 5 électrons de valence. Trois d'entre eux servent aux 3 liaisons \ce{N-H} ; il
        reste $5-3 = 2$ électrons, soit \textbf{1 doublet non liant}.
  \item L'azote est le \textbf{donneur} : il apporte \emph{les deux} électrons de son doublet
        libre. Le proton \ce{H+}, qui possède une case vide et aucun électron, est l'
        \textbf{accepteur}. Le doublet est mis en commun : c'est une liaison covalente de
        coordination (dative).
  \item Une fois la liaison formée, rien ne distingue physiquement le doublet « donné » d'un
        doublet ordinaire : les électrons appartiennent aux deux atomes. Les quatre liaisons
        \ce{N-H} sont donc \textbf{équivalentes} et l'ion \ce{NH4+} est parfaitement symétrique
        (tétraédrique).
  \item L'azote forme 4 liaisons et ne porte plus aucun doublet libre : $4\times 2 = \mathbf{8}$
        électrons. \rep l'octet est \textbf{respecté}.
\end{enumerate}
\end{cor}

\begin{cor}[trois atomes centraux, trois comportements face à l'octet]
\begin{enumerate}[label=\alph*),leftmargin=2em,itemsep=2pt]
  \item \ce{BF3} : $3\times 2 = 6$ électrons. \quad \ce{CF4} : $4\times 2 = 8$. \quad
        \ce{SF6} : $6\times 2 = 12$.
  \item \ce{BF3} : octet \textbf{incomplet} (6). \ce{CF4} : octet \textbf{respecté} (8).
        \ce{SF6} : octet \textbf{étendu} (12).
  \item Pour \ce{BF3} : le \textbf{bore} n'a que 3 électrons de valence ; il se contente d'un
        « sextet », c'est une exception admise (il n'a pas assez d'électrons pour un octet sans
        charges). Pour \ce{SF6} : le \textbf{soufre} est en \textbf{3\textsuperscript{e} période}
        et possède des orbitales $d$ vides ; il peut donc loger plus de 8 électrons (hypervalence).
  \item \textbf{Non}, « \ce{OF6} » est impossible : l'\textbf{oxygène} est en 2\textsuperscript{e}
        période, il ne dispose d'aucune orbitale $d$ et ne peut \textbf{jamais dépasser 8}
        électrons. Six liaisons (12 électrons) sont donc exclues.
\end{enumerate}
\end{cor}

\begin{cor}[le cristal ionique : vrai ou faux]
\begin{enumerate}[label=\alph*),leftmargin=2em,itemsep=2pt]
  \item \textbf{Faux.} La cohésion est due aux interactions \textbf{électrostatiques
        (coulombiennes)} entre ions de charges opposées, pas aux interactions gravitationnelles
        (infimes à l'échelle atomique).
  \item \textbf{Vrai.} Le cristal est un empilement \emph{ordonné} et régulier d'ions ; la
        dissolution dans l'eau \textbf{sépare et disperse} ces ions (solvatés par les molécules
        d'eau), ce qui détruit l'ordre cristallin.
  \item \textbf{Vrai.} La solution contient des \textbf{ions mobiles} (\ce{Na+} et \ce{Cl-}
        libres) : c'est une solution \textbf{électrolytique}, qui conduit le courant.
\end{enumerate}
\end{cor}

\begin{astucebox}[la logique commune aux trois types de liaison]
Tout part du même moteur : chaque atome cherche à compléter sa couche externe (octet/duet).
Il peut le faire en \textbf{partageant} (covalente), en \textbf{transférant} (ionique), ou en
\textbf{recevant un doublet tout fait} (dative). Identifier le mécanisme, c'est déjà comprendre la
molécule.
\end{astucebox}

\end{document}
